\par\medskip\noindent
\underline{Proposición} 1.3.5. \underline{Sean} $P$ \underline{y}
$\sisi{Q}{G}$ \underline{dos grupos de} $\underline{T}$.

a)  \underline{Supongamos que todo morfismo de haces de conjuntos de}
$P\times P$ \underline{en} $G$ \underline{que sea trivial sobre}
$P\times e_P$ \underline{y sobre} $e_P\times P$ \underline{es
trivial}. \underline{Entonces el funtor} $i$ \underline{de} (1.3.4.1)
\underline{es plenamente fiel}. \underline{En particular},
\underline{sobre un bitorsor} $E$ \underline{bajo} $(G_P,G_P)$,
\underline{\sisi{birrigidificado}{rigidificado} relativamente a} $e_P$, \underline{es
decir}, \underline{provisto de una sección central} $e_E$
\underline{de} $E_1 = e_P^*(E)$, \underline{existe a lo sumo una
estructura de extensión de} $P$ \underline{por} $G$,
\underline{compatible con su estructura de bitorsor y que admite}
$e_E$ \underline{como sección unidad}.

b)  \underline{Supongamos además que todo morfismo de haces de
conjuntos de} $P\times P\times P$ \underline{en} $G$ \underline{que
sea trivial sobre los productos parciales} $P\times P\times e_P$,
$P\times e_P\times P$ \underline{y} $e_P\times P\times P$
\underline{es trivial}. \underline{Sea} $E$
\underline{\sisi{un torsor bajo $G_P$}{un bitorsor bajo
$(G_P,G_P)$}}. \underline{Para que exista sobre} $E$ \underline{una
estructura de extensión de} $P$ \underline{por} $G$
\underline{compatible con su estructura de bitorsor}, \underline{es
necesario y suficiente que el bitorsor bajo}
$(G_{P\times P},G_{P\times P})$ \underline{subyacente a}
$\delta(E) = \pi^*(E)\mathrm{pr}_1^*(E)^{-1}\mathrm{pr}_2^*(E)^{-1}$
\underline{sea trivial}. \underline{Cuando así sea}, \underline{para
toda sección central} $e_E$ de $E_1 = e_P^*(E)$, \underline{es decir},
\underline{toda rigidificación de} $E$ \underline{relativamente a}
$e_P$, \underline{existe una estructura de extensión única sobre}
$E$, \underline{compatible con su estructura de bitorsor y que admite}
$e_E$ \underline{como sección unidad}; \underline{es decir}, $E$,
\underline{considerado como bitorsor rigidificado relativamente a}
$e_P$, \underline{pertenece a la imagen esencial de} $i$. \underline{Por
otra parte}, \underline{existe una y solo una trivialización del
bitorsor birrigidificado} $\delta(E)$.


%%%%%%%%%%  scan idx 160  |  folio 149  %%%%%%%%%%

En términos del diagrama (1.3.4.1), puede decirse, pues, que la imagen
esencial del funtor plenamente fiel $i$ es «el núcleo» del funtor
$\delta$, o también, de manera figurada, que el diagrama (1.3.4.1) es
una «sucesión exacta de categorías apuntadas», e incluso de categorías
con estructuras de pseudogrupos (en la terminología introducida en
2.5.1 más adelante).

Demostremos 1.3.5.  a)  Sean $E,E'$ dos extensiones de $P$ por $G$, y
sea $f : E \longrightarrow E'$ un morfismo de los bitorsors
rigidificados subyacentes; demostremos que $f$ es un morfismo de
extensiones o, lo que manifiestamente equivale a lo mismo, un morfismo
de grupos. Hay que demostrar que $f(xy) = f(x)f(y)$. Ahora bien
(designando por $j' : E' \longrightarrow P$ el morfismo estructural del
torsor $E'$ sobre $P$), se tendrá
$j'f(xy) = j'(f(x)f(y))$, pues los dos miembros son respectivamente
iguales a $\sisi{j(f(xy))}{j(xy)}$ y a $\sisi{j(f(x)f(y))}{j(x)j(y)}$, que son iguales porque $j$ es
multiplicativo. Por tanto,

\begin{equation*}
f(xy) = u(j(x),j(y))f(x)f(y) \qquad ,
\end{equation*}
donde

\begin{equation*}
u : P \times P \longrightarrow G
\end{equation*}
es un morfismo bien determinado. Haciendo $x=1$ y después $y=1$ en la
relación precedente, se obtiene (teniendo en cuenta que $f(1)=1$)

\begin{equation*}
u(p,1) = u(1,p) = 1 \qquad ,
\end{equation*}
es decir, $u$ es trivial sobre $P \times e_P$ y sobre
$e_P \times P$, de donde resulta por hipótesis que $u$ es trivial, y
por tanto $f(xy) = f(x)f(y)$, como se quería demostrar.

b)  La necesidad de la condición considerada sobre $E$ para que
admita una estructura de extensión compatible con su estructura de
bitorsor es trivial (1.3.4.3); queda demostrar su suficiencia para
probar la primera afirmación de b). Ahora bien, utilizando una
trivialización $\varphi$
%%%%%%%%%%  scan idx 161  |  folio 150  %%%%%%%%%%
% Canon: SGA-SRC-001257.
de $\delta(E)$, se deduce de ella, por imagen inversa mediante la
sección unidad de $P \times P$, una trivialización de
$E_1 = \sisi{e_P(E)}{e_P^*(E)}$ (compárese 1.5.1), es decir, una
rigidificación de $E$
relativamente a $e_P$. Supongámosla fijada. Una trivialización dada
$\varphi$ de $\delta(E)$ no es necesariamente compatible con esta
birrigidificación; sus restricciones a $P \times e_P$ y a
$e_P \times P$ se expresan, en términos de las rigidificaciones
parciales naturales de $\delta(E)$, mediante morfismos
\underline{centrales}

\begin{equation*}
f : P \longrightarrow G \ , \ g : P \longrightarrow G \ ,
\end{equation*}
que tienen las mismas restricciones a la sección unidad de $P$:
$f(1) = g(1) = a \in \Gamma G$. Entonces el morfismo
\underline{central}

\begin{equation*}
h : P \times P \longrightarrow G \quad , \quad h(p,q) = f(p)g(q)a^{-1}
\end{equation*}
tiene como restricciones a $P \times e_P$ y $e_P \times P$ los
morfismos $f$ y $g$. Por consiguiente, corrigiendo $\varphi$ mediante
$h$ (es decir, reemplazándolo por $\varphi(p,q)h(p,q)^{-1}$), se
obtiene una trivialización de $\delta(E)$ compatible con su
birrigidificación. Observemos, por lo demás, que una trivialización de
esta índole está determinada módulo un morfismo central
$P \times P \longrightarrow G$ que es trivial sobre $P \times e_P$
y sobre $e_P \times P$, y que por tanto es trivial; así, este
$\varphi$ es único, lo que constituye la última afirmación de 1.3.5 b).
La demostración quedará terminada si probamos que este $\varphi$,
considerado como un isomorfismo de la forma (1.1.3.2), satisface la
condición de asociatividad (1.1.4.1), teniendo en cuenta que ya sabemos,
por la elección de $\varphi$, que $e_E$ es una unidad bilateral para la
ley de composición definida por $\varphi$. Ahora bien, los dos
compuestos de (1.1.4.1) son isomorfismos de bitorsors bajo
$(G_{P \times P \times P} , G_{P \times P \times P})$, por lo que
difieren en un morfismo central

\begin{equation*}
u : P \times P \times P \longrightarrow G \qquad ,
\end{equation*}
según la fórmula

\begin{equation*}
(xy)z = u(jx,jy,jz) \ x(yz) \qquad .
\end{equation*}

%%%%%%%%%%  scan idx 162  |  folio 151  %%%%%%%%%%

Haciendo en esta fórmula sucesivamente
\sisi{$y=z=1$, después $z=x=1$, y después $x=y=1$}{$z=1$, después
$x=1$, y después $y=1$}, y utilizando el hecho de que $1$ es una
unidad bilateral para la ley de composición considerada, se obtiene que
$u$ es trivial sobre los tres trozos
\sisi{$P\times e_P\times e_P$, $e_P\times P\times e_P$,
$e_P\times e_P\times P$}{$P\times P\times e_P$,
$e_P\times P\times P$, $P\times e_P\times P$}, lo cual implica por
hipótesis que es trivial, y por tanto que $(xy)z = x(yz)$, con lo que
termina la demostración.
